\documentclass[11pt]{article}
\usepackage{amsmath,amssymb,geometry,hyperref}
\geometry{margin=1in}
\title{From-Scratch Replication: Magnetostriction as a Probe of\\
Octupolar Order under a [111] Field\\
\large (Patri \textit{et al.}, arXiv:1901.00012, 2018)}
\author{TEXTURES-100 replication (Hermes subagent)}
\date{\today}
\begin{document}
\maketitle

\section{Model and System}
The paper treats the Pr-based cage compounds Pr(Ti,V,Ir)$_2$(Al,Zn)$_{20}$
(illustratively PrV$_2$Al$_{20}$). The Pr$^{3+}$ $4f^2$ ions on the diamond
lattice (space group $Fd\bar{3}m$, local $T_d$ symmetry) host a non-Kramers
$\Gamma_{3g}$ ground doublet, well separated ($\gtrsim 50$~K) from excited
states. The doublet carries \emph{no} magnetic dipole moment; its only low-energy
degrees of freedom are the quadrupoles $O_{20},O_{22}$ (time-reversal even) and
the octupole $T_{xyz}=\tfrac{\sqrt{15}}{6}\overline{J_xJ_yJ_z}$ (time-reversal
odd, $\Gamma_{2u}$). These map onto a pseudospin-$\tfrac12$:
\begin{equation}
\tau^x=-\tfrac14 O_{22},\quad \tau^y=-\tfrac14 O_{20},\quad
\tau^z=\tfrac{1}{3\sqrt5}T_{xyz}.
\end{equation}
Order parameters: AFQ $\tilde\varphi$, FQ $\varphi$ (parasitic), and the
ferro-octupole (FO) $m\equiv\langle\tau^z_A+\tau^z_B\rangle$.

\section{Method (built from scratch)}
Symmetry-based Landau theory plus a strain minimization, reusing the shared
kernel \texttt{ollie\_multipolar\_stevens\_landau\_kernel.py} (Stevens operators,
thermal susceptibility) for the single-ion multipole construction (provenance
credit). Three independent computations:
\begin{enumerate}
\item \textbf{(A) Pseudospin construction.} Rebuild the $\Gamma_3$ doublet from
the $J=4$ multiplet (Eq.~2 of paper), form $|\!\uparrow\rangle,|\!\downarrow\rangle$
(Eq.~4), and verify $\tau^z\propto T_{xyz}$ maps onto a genuine pseudospin-$\tfrac12$
operator. Also compute the single-ion octupole susceptibility $\chi_{T_{xyz}}(T)$.
\item \textbf{(B) Strain minimization (the headline).} With the cubic elastic
energy
\begin{equation}
F_{\rm lat}=\tfrac{c_{11}}{2}(\varepsilon_{xx}^2{+}\varepsilon_{yy}^2{+}\varepsilon_{zz}^2)
+\tfrac{c_{44}}{2}(\varepsilon_{xy}^2{+}\varepsilon_{yz}^2{+}\varepsilon_{xz}^2)
+c_{12}(\varepsilon_{xx}\varepsilon_{yy}{+}\varepsilon_{yy}\varepsilon_{zz}{+}\varepsilon_{zz}\varepsilon_{xx})
\end{equation}
and the field-induced octupole--strain coupling
$\Delta F=-g_O\,m(\varepsilon_{yz}h_x+\varepsilon_{xz}h_y+\varepsilon_{xy}h_z)$,
minimize $F_{\rm lat}+\Delta F$ over the full strain tensor. This gives
$\varepsilon_{xy}=(g_O/c_{44})m\,h_z$ (cyclically), diagonal strains vanish, and
for $\mathbf{h}\parallel[111]$ ($h_i=h/\sqrt3$),
\begin{equation}
\left(\frac{\Delta L}{L}\right)_{[111]}=\frac{\varepsilon_{xy}+\varepsilon_{yz}+\varepsilon_{xz}}{3}
=\frac{1}{\sqrt3}\,\frac{g_O}{c_{44}}\,m\,h .
\end{equation}
\item \textbf{(C) Hysteresis.} Minimize the FO Landau potential
$F_m=\tfrac{t_m}{2}m^2+u_m m^4-b\,(h_xh_yh_z)\,m$ (the cubic-in-$h$, $\sim h_xh_yh_z\tau^z$
drive) on up/down field sweeps below $T_O$ ($t_m<0$) and measure the loop width.
\end{enumerate}

\section{Results}
\begin{itemize}
\item \textbf{(A)} Doublet orthonormal (overlaps $10^{-16}$). $\tau^z$ from
$T_{xyz}$ matches $\tfrac12\sigma_z$ with residual $4\times10^{-17}$; $\tau^x$ from
$O_{22}$ residual $6\times10^{-17}$. Octupole susceptibility is Curie-like:
$\chi_{T_{xyz}}=220$ at $T=0.3$ vs $22$ at $T=3.0$ (ratio exactly $1/T$),
supporting a ferro-octupolar instability at low $T$.
\item \textbf{(B)} Log--log exponent of $(\Delta L/L)_{[111]}$ vs $h$ is
$1.0000$ (\emph{exactly linear in $h$}). Coefficient vs $m$: $R^2=1.000$
(\emph{exactly $\propto m$}). Coefficient magnitude
$A/[(g_O/c_{44})m]=0.57735=1/\sqrt3$ --- exactly the [111] geometric projection.
\item \textbf{(C)} Below $T_O$ the octupole $m(h)$ traces a hysteresis loop of
width $1.00$ (in units of the spontaneous $m$), with a finite coercive field; the
$[111]$ length change inherits this hysteresis.
\end{itemize}

\section{Comparison to the Claim}
\begin{center}
\begin{tabular}{lll}
\hline
Claim & Paper & This work\\
\hline
$(\Delta L/L)_{[111]}\sim h^1$ & linear-in-$h$ & exponent $=1.0000$ \checkmark\\
coeff.\ $\propto m$ (FO) & yes & $R^2=1.000$ \checkmark\\
coeff.\ $\propto g_O/c_{44}$ & yes & ratio $=1/\sqrt3$ (geom.) \checkmark\\
hysteresis below $T_O$ & yes & loop width $=1.0$ \checkmark\\
\hline
\end{tabular}
\end{center}
All four sub-claims of the headline are reproduced analytically/numerically.
\textbf{Verdict: REPLICATED.}

\section{Critique}
The result is a symmetry-dictated identity rather than a material-specific
number: given the $T_d$-allowed couplings, the linear-in-$h$ scaling and the
$\propto m$ proportionality follow \emph{exactly} and are essentially
parameter-free (only the geometric $1/\sqrt3$ enters for [111]). Thus the
replication confirms the paper's \emph{internal logic} strongly, but does not
independently validate (i) the microscopic value of $g_O$ (needs the second/third
order perturbation theory in $\mathbf{h}\!\cdot\!\mathbf{J}$ and CEF gaps
$\Delta(\Gamma_4),\Delta(\Gamma_5)$), nor (ii) the actual value of $T_O\approx
0.65$~K or the loop shape in PrV$_2$Al$_{20}$, which require the full
temperature-dependent coupled minimization of $F_{\tilde\varphi,\varphi,m}$ and
experimental $c_{44}$, $g_O$. Our hysteresis is a demonstrator (single-well FO
potential with cubic drive), not a fit to data. See
\texttt{failure\_analysis.md}.

\section{Provenance}
Single-ion angular-momentum matrices, Stevens multipole operators
($O_{20},O_{22},T_{xyz}$), and the thermal (fluctuation) susceptibility were
reused from \texttt{ollie\_multipolar\_stevens\_landau\_kernel.py} (Ollie's shared
TEXTURES-100 multipolar kernel). All Landau/strain minimization and scaling
analysis are original to this replication.

\end{document}
